\subsection{What finite quadratic evaluations can and cannot detect}

This subsection records a bounded attempt to detect the explicit
Eisenstein class in KT-EISENSTEIN-QUADRATIC-DETECTOR.  It proves two
limitations of a proposed spanning argument and a sufficient short
discriminant window.  It does not prove that this window contains the
required field, and it gives no additional vanishing of a rational
class-group component.

Fix a prime $p\geq31$ with $p\equiv3\pmod4$, and put
\[
 b=(p+3)/2,\qquad k=(p+9)/2,\qquad e=b-1.
\]
For the actual auxiliary-field problem the three units are
$B_{1,\chi_D}$, $\beta_b(\chi_D)$ and $\beta_3(\chi_D)$:
the last index is obtained from $2e+1=p+2$ by the previously
checked Kummer transport.  Here $\beta_r(\chi)=B_{r,\chi}/r$.
The discriminant must remain odd and negative and prime to $p$.
The carried small-discriminant bound gives $p\nmid h(D)$ whenever
$|D|<4p$.

\begin{proposition}[KT-SMALL-QUADRATIC-SPAN-LIMIT]
Let $q$ be prime with $q\equiv1\pmod p$.  Consider the selected
$j=2$ quadratic evaluation functionals at the odd negative fundamental
discriminants $D$ satisfying
\[
 0<|D|<4p,\qquad \gcd(D,pq)=1,\qquad \chi_D(q)=-1.
\]
On a full integral lattice in either parity component of the
characteristic-zero cuspidal modular-symbol space of weight $k$ and
level $\Gamma_0(q)$, these functionals cannot span the entire dual
lattice, either rationally or after reduction modulo $p$ when the
functionals are integral on that lattice.
\end{proposition}

\begin{proof}
There are at most $p$ candidate integers: an odd negative fundamental
discriminant is $D=-d$ with $d\equiv3\pmod4$, and there are $p$ such
positive integers below $4p$.  Squarefreeness and the displayed local
conditions only decrease their number.

Either cuspidal parity component has rank
$s=\dim S_k(\Gamma_0(q))$.  In the standard dimension formula put
$e_2=1+(\frac{-4}{q})$ and $e_3=1+(\frac{-3}{q})$, so that
$0\leq e_2,e_3\leq2$.  Since there are two cusps, the formula is
\[
 s=\frac{(k-1)(q+1)}{12}-1
   +e_2\left(\left\lfloor\frac{k}{4}\right\rfloor
                         -\frac{k-1}{4}\right)
   +e_3\left(\left\lfloor\frac{k}{3}\right\rfloor
                         -\frac{k-1}{3}\right).
\]
For even $k$, the two parenthetical terms are at least $-1/4$
and $-1/3$, respectively.  Also $q\geq2p+1$, because both primes
are odd.  Hence
\[
 s\geq\frac{(p+7)(p+1)}{12}-\frac{13}{6}>p.
\]
The last difference is $(p^2-4p-19)/12$, positive for $p\geq31$.
At most $p$ vectors cannot span a space of dimension greater than
$p$.  Reduction of a free lattice still has its full rank as
dimension over $\mathbf F_p$, so the same conclusion holds there.
\end{proof}

This proposition does not say that the explicit Eisenstein class is
annihilated by the small-discriminant evaluations.  Nor does it bound
an Eisenstein-localized eigenspace by $s$.  It rules out the proposed
full-lattice spanning shortcut, not a proof of detection of that
particular class.

\begin{lemma}[KT-EISENSTEIN-HECKE-EVALUATION-IDEAL]
Let $R$ be a ring, let $M$ be an $R$-module with Hecke action, and
let $v\in M$ satisfy $Tv=a(T)v$ for every operator being used.
For any family of $R$-linear functionals $\lambda_D$, the ideal
\[
 (\lambda_D(Tv):D,T)\subset R
\]
is exactly $(\lambda_D(v):D)$ if the identity operator is included.
\end{lemma}

\begin{proof}
One has $\lambda_D(Tv)=a(T)\lambda_D(v)$, giving one inclusion.
The identity operator gives the reverse inclusion.
\end{proof}

For the Heumann--Vatsal eigenclass and the chosen local condition,
the generators in this lemma are the unit multiple
$-\frac{15}{8}\beta_3(\chi_D)\beta_b(\chi_D)$.
Thus taking the Hecke closure of the evaluation rows cannot itself
improve their scalar ideal on this class.  This observation does
not rule out proving that the original ideal is the unit ideal.
Varying $q$ likewise changes only the already displayed Euler
factor.  For this pair of Bernoulli indices, Corollary 4.7 has
$j=2$ or $j=k-4$; an unspecified nonzero critical coefficient
can concern different Bernoulli factors.

\paragraph{The exact finite quadratic orthogonality calculation.}
Let $m$ be odd and squarefree, prime to $pq$, with $m<4p$.
Put $G=(\mathbf Z/m\mathbf Z)^\times$ and
$X=\operatorname{Hom}(G,\{\pm1\})$.  For $r=3,b$ define the
$\mathbf F_p$-valued function
\[
 F_r(a)=\frac{m^{r-1}}r B_r(a/m)\pmod p,\qquad a\in G.
\]
All denominators are units: $r<p-1$ and $p\nmid m$.
The Bernoulli-polynomial definition gives
\[
 \beta_{r,m}(\chi)=\sum_{a\in G}\chi(a)F_r(a).
\]
The subscript $m$ allows the character to be imprimitive.
Writing $H=G^2$, ordinary finite-group character orthogonality gives
\[
 \sum_{\substack{\chi\in X\\\chi(-1)=\chi(q)=-1}}\chi(g)
 =\frac{|X|}{4}
  \bigl(1_H(g)-1_H(-g)-1_H(qg)+1_H(-qg)\bigr).
 \tag{QE.1}
\]
Indeed multiply the unrestricted sum by
$(1-\chi(-1))(1-\chi(q))/4$ before summing; this proof also
covers dependent local conditions or an empty selected family.
Consequently the sum of the products
$\beta_{3,m}(\chi)\beta_{b,m}(\chi)$ over the selected characters
is exactly
\[
 \frac{|X|}{4}\sum_{a,c\in G}F_3(a)F_b(c)
  \bigl(1_H(ac)-1_H(-ac)-1_H(qac)+1_H(-qac)\bigr).
 \tag{QE.2}
\]
This is a concrete finite expression, with the same character in
both factors.  A nonzero value proves that one product is a unit.
If $\chi^*$ is the primitive character inducing $\chi$, then
\[
 \beta_{r,m}(\chi)=\beta_r(\chi^*)
  \prod_{\substack{\ell\mid m\\\ell\nmid\operatorname{cond}(\chi^*)}}
       (1-\chi^*(\ell)\ell^{r-1}).
 \tag{QE.3}
\]
Thus a nonzero product in (QE.2) supplies the required two primitive
units and a discriminant of size at most $m<4p$; its class-number
unit follows from the carried bound.  The parity and local
conditions also pass to $\chi^*$.

No nonzero value of (QE.2) has been proved uniformly here.
Quadratic orthogonality detects only the sums of $F_r$ on the
squareclasses $G/H$, not its individual values.  In particular
$\delta_a-\delta_{ah}$ is in its kernel for $h\in H$ with
$ah\ne a$.  Even nonzero quadratic Fourier transforms of two
functions may have disjoint supports.  Therefore separate
nonvanishing of the two transforms does not prove that their
pointwise product is nonzero.  These are exact limitations of an
argument using only character orthogonality, not counterexamples
to nonvanishing for the specific Bernoulli functions $F_r$.
The factors in (QE.3) can also vanish modulo $p$, so a primitive
unit need not be visible in a chosen imprimitive conductor.

\begin{proposition}[KT-BETA3-SHORT-WINDOW]
Let $D=-d$ be an odd negative fundamental discriminant.
If $d\ne3,7$, then $\beta_3(\chi_D)$ is a positive integer and
\[
 0<\beta_3(\chi_D)<\frac{5}{216}d^{5/2}.
 \tag{QE.4}
\]
For $d=3,7$ the exact values are $2/9$ and $16/7$.
In particular, for every prime $p\geq31$, the condition
\[
 p\nmid d,\qquad d^5\leq1600p^2
 \tag{QE.5}
\]
implies that $\beta_3(\chi_D)$ is a $p$-adic unit.
\end{proposition}

\begin{proof}
Write $\chi=\chi_D=(\frac{\cdot}{d})$ and
$S_j=\sum_{a=1}^d\chi(a)a^j$.  Odd parity, paired by $a\mapsto d-a$,
gives $S_2=dS_1$, and the explicit polynomial
$B_3(x)=x^3-\frac32x^2+\frac12x$ gives
\[
 \beta_3(\chi)=\frac{S_3-d^2S_1}{3d}.
 \tag{QE.6}
\]
For each prime $\ell\mid d$, the Chinese remainder theorem shows
$S_3\equiv0\pmod\ell$ if $d/\ell>1$: the sum of the nontrivial
quadratic character modulo $d/\ell$ is zero.  If $d=\ell$ is prime,
Euler's criterion and the finite-field power sum give the same
conclusion unless $\ell=3$ or $7$.  Indeed the relevant exponent
is $(\ell-1)/2+3$, which is not a multiple of $\ell-1$ for
$\ell>7$.  Thus $d\mid S_3$ outside the two stated cases.

If $3\nmid d$, then $a^3\equiv a\pmod3$ and $d^2\equiv1\pmod3$,
so the numerator in (QE.6) is also divisible by $3$.
If $3\mid d$ and $d>3$, put $d=3m$.  For $3\nmid a$,
$(\frac a3)a^3\equiv1\pmod9$.  Hence
\[
 S_3\equiv
 \sum_{\substack{1\leq a\leq3m\\3\nmid a}}
          \left(\frac a m\right)=0\pmod9.
\]
The last equality follows by subtracting the sum over multiples
of $3$ from three complete periods of the nontrivial character
modulo $m$.  Also $9\mid d^2S_1$.  Since $d$ is squarefree, these
facts prove integrality in all the required cases.
Direct substitution in (QE.6) gives the two exceptional values.

The same Bernoulli Fourier calculation as in (QD.1), using
$\tau(\chi_D)=i\sqrt d$, gives
\[
 \beta_3(\chi_D)=\frac{d^{5/2}}{2\pi^3}L(3,\chi_D).
\]
The Euler product is positive and
$L(3,\chi_D)\leq\zeta(3)<5/4$; the latter follows from
$1+2^{-3}+\int_2^\infty x^{-3}\,dx=5/4$.
Together with $\pi>3$ this proves (QE.4).
Under (QE.5), its upper bound is at most
$200p/216=25p/27<p$.  A positive integer smaller than $p$
is a unit modulo $p$.  The exceptional values are also units
for $p\geq31$.
\end{proof}

Therefore a $\beta_b$-unit seed in the window (QE.5) would supply
the same-discriminant product and the class-number condition.
This window has size $O(p^{2/5})$; the established Jacobi
heat/Sturm construction gives a bound of order $p$ and has not
been improved to this window.  No existence assertion in (QE.5)
is made.

\paragraph{Source and verification scope.}
The modular-symbol and Bernoulli normalization is the already
whole-reviewed Heumann--Vatsal Corollary 4.7 and
KT-EISENSTEIN-QUADRATIC-DETECTOR.  The dimension input is
W.~Stein, \emph{Modular Forms: A Computational Approach},
Proposition 6.1, printed p.~93, with its definitions on p.~92,
and the parity pairing in Theorem 8.23, printed p.~141,
in the primary author PDF
\url{https://wstein.org/books/modform/stein-modform.pdf}.
Those precise statements were read; their foundational
Riemann--Roch and period-pairing inputs are imported,
not re-proved here.
The elementary character, scalar-ideal, and integrality
arguments above are given in full.  This is a limitation and
sufficient-window result, not arithmetic movement on the
remaining rational component.
